% !TEX encoding = UTF-8 Unicode
\documentclass[11pt,a4paper]{article}

\usepackage[utf8]{inputenc}
\usepackage[T1]{fontenc}
\usepackage{geometry}
\usepackage{hyperref}
\usepackage{enumitem}
\usepackage{booktabs}
\usepackage{array}
\usepackage{tabularx}
\usepackage{xcolor}
\usepackage{fancyhdr}
\usepackage{titlesec}
\usepackage{parskip}
\usepackage{listings}
\usepackage{tcolorbox}
\tcbuselibrary{breakable, skins, listings}

\geometry{a4paper, top=2.2cm, bottom=2.2cm, left=2.5cm, right=2.5cm}

\definecolor{accentblue}{RGB}{30, 80, 160}
\definecolor{promptbg}{RGB}{242, 244, 248}
\definecolor{promptborder}{RGB}{30, 80, 160}
\definecolor{warnbg}{RGB}{255, 248, 235}
\definecolor{warnborder}{RGB}{190, 80, 0}
\definecolor{tipbg}{RGB}{240, 248, 240}
\definecolor{tipborder}{RGB}{50, 130, 60}
\definecolor{darkgray}{RGB}{70, 70, 70}
\definecolor{medgray}{RGB}{130, 130, 130}
\definecolor{cheatbg}{RGB}{250, 250, 252}
\definecolor{cheatborder}{RGB}{30, 80, 160}
\definecolor{labelcolor}{RGB}{40, 40, 100}
\definecolor{notebg}{RGB}{245, 245, 255}
\definecolor{noteborder}{RGB}{100, 100, 180}

\pagestyle{fancy}
\fancyhf{}
\fancyhead[L]{\small\textcolor{medgray}{AI-alapú részvényelemzés}}
\fancyhead[R]{\small\textcolor{medgray}{Sapientia EMTE -- dr.\ Pál László}}
\fancyfoot[C]{\small\textcolor{medgray}{\thepage}}
\renewcommand{\headrulewidth}{0.3pt}

\titleformat{\section}{\large\bfseries\color{accentblue}}{\thesection}{0.6em}{}%
  [\vspace{-2pt}\textcolor{accentblue}{\rule{\linewidth}{0.6pt}}]
\setcounter{secnumdepth}{0}
\titleformat{\subsection}{\normalsize\bfseries\color{darkgray}}{}{0em}{}
\titlespacing{\section}{0pt}{18pt}{8pt}
\titlespacing{\subsection}{0pt}{12pt}{4pt}

\lstdefinestyle{promptlst}{
  basicstyle=\small\ttfamily\color{black},
  breaklines=true,
  breakatwhitespace=false,
  frame=none,
  columns=fullflexible,
  keepspaces=true,
  showstringspaces=false,
  literate=
    {á}{{\'{a}}}1 {é}{{\'{e}}}1 {í}{{\'{i}}}1 {ó}{{\'{o}}}1 {ö}{{\"o}}1
    {ő}{{\H{o}}}1 {ú}{{\'{u}}}1 {ü}{{\"u}}1 {ű}{{\H{u}}}1
    {Á}{{\'{A}}}1 {É}{{\'{E}}}1 {Í}{{\'{I}}}1 {Ó}{{\'{O}}}1 {Ö}{{\"O}}1
    {Ő}{{\H{O}}}1 {Ú}{{\'{U}}}1 {Ü}{{\"U}}1 {Ű}{{\H{U}}}1,
}

\tcbset{
  promptstyle/.style={
    breakable, enhanced,
    colback=promptbg, colframe=promptborder,
    boxrule=1pt, arc=3pt,
    left=10pt, right=10pt, top=10pt, bottom=8pt,
    attach boxed title to top left={xshift=10pt, yshift=-\tcboxedtitleheight/2},
    boxed title style={
      colback=promptborder, colframe=promptborder,
      arc=2pt, boxrule=0pt,
      left=6pt, right=6pt, top=2pt, bottom=2pt
    },
    coltitle=white, fonttitle=\footnotesize\bfseries\sffamily,
  },
  warnstyle/.style={
    breakable, colback=warnbg, colframe=warnborder,
    boxrule=0.5pt, leftrule=3pt, arc=2pt,
    left=8pt, right=8pt, top=5pt, bottom=5pt,
    fontupper=\small\sffamily,
  },
  tipstyle/.style={
    breakable, colback=tipbg, colframe=tipborder,
    boxrule=0.5pt, leftrule=3pt, arc=2pt,
    left=8pt, right=8pt, top=5pt, bottom=5pt,
    fontupper=\small\sffamily,
  },
  cheatstyle/.style={
    breakable, enhanced, colback=cheatbg, colframe=cheatborder,
    boxrule=0.8pt, arc=3pt,
    left=10pt, right=10pt, top=9pt, bottom=7pt,
    attach boxed title to top left={xshift=10pt, yshift=-\tcboxedtitleheight/2},
    boxed title style={
      colback=cheatborder, colframe=cheatborder,
      arc=2pt, boxrule=0pt,
      left=6pt, right=6pt, top=2pt, bottom=2pt
    },
    coltitle=white, fonttitle=\footnotesize\bfseries\sffamily,
  },
  notestyle/.style={
    breakable, colback=notebg, colframe=noteborder,
    boxrule=0.5pt, leftrule=3pt, arc=2pt,
    left=8pt, right=8pt, top=5pt, bottom=5pt,
    fontupper=\small\sffamily,
  },
}

\newtcblisting{prompt}[1][Prompt]{
  promptstyle, title=#1,
  listing only,
  listing options={style=promptlst},
}

\newcommand{\warn}[1]{%
  \vspace{3pt}%
  \begin{tcolorbox}[warnstyle]%
  \textbf{\textcolor{warnborder}{Figyelem:}} #1%
  \end{tcolorbox}\vspace{2pt}}

\newcommand{\tip}[1]{%
  \vspace{3pt}%
  \begin{tcolorbox}[tipstyle]%
  \textbf{\textcolor{tipborder}{Tipp:}} #1%
  \end{tcolorbox}\vspace{2pt}}

\newcommand{\note}[1]{%
  \vspace{3pt}%
  \begin{tcolorbox}[notestyle]%
  \textbf{\textcolor{noteborder}{Változtatási javaslat:}} #1%
  \end{tcolorbox}\vspace{2pt}}

\newcommand{\lbl}[1]{\vspace{6pt}\noindent{\small\bfseries\textcolor{labelcolor}{#1}}\par\smallskip}

\hypersetup{colorlinks=true, linkcolor=accentblue, urlcolor=accentblue}

\begin{document}

\thispagestyle{empty}
\begin{center}
  {\LARGE\bfseries AI-alapú részvényelemzés}\\[5pt]
  {\normalsize Generatív AI és alkalmazások modul}\\[1pt]
  {\small\color{medgray} dr.\ Pál László, Sapientia EMTE, 2026. március 6--7.}
\end{center}

\vspace{10pt}
\noindent
\begin{tabularx}{\linewidth}{lXr}
\toprule
\textbf{Blokk} & \textbf{Tartalom} & \textbf{Időtartam} \\
\midrule
1.\ rész & Adatletöltés + alapvizualizáció (árfolyam, napi hozam, kumulált hozam) & $\sim$30 perc \\
2.\ rész & FinBERT sentiment elemzés & $\sim$40 perc \\
3.\ rész & Kereskedési stratégia (technikai indikátorok + backtesting) & $\sim$30 perc \\
\midrule
\multicolumn{2}{l}{\textbf{Teljes blokk}} & \textbf{$\sim$100 perc} \\
\bottomrule
\end{tabularx}

\smallskip
\noindent\textbf{Eszközök:} Google Colab (\texttt{colab.research.google.com}),
Gemini (\texttt{gemini.google.com})

\vspace{8pt}

%=============================================================
\section{1. RÉSZ -- Adatletöltés és alapvizualizáció}
%=============================================================

\noindent\textit{Cél: Letöltjük a részvény és az S\&P~500 index historikus adatait.
Az árfolyamot USD-ben vizualizáljuk, a napi hozamot csak a részvényre nézzük,
a kumulált hozamnál pedig megjelenik az index összehasonlítás.
Az első két cella (telepítés, adatletöltés) előre megírva szerepel a notebookban -- csak futtatni kell.}

\bigskip

%--- Telepítés ---
\subsection{0. lépés -- Telepítés és adatletöltés}

\noindent Futtasd le sorban a notebook első celláit:

\begin{enumerate}[noitemsep, topsep=4pt]
  \item \textbf{Telepítés / importok} -- csomagok betöltése
  \item \textbf{Adatletöltés} -- az AAPL részvény \textit{és} az S\&P~500 index (\texttt{\^{}GSPC}) letöltése 2020-tól 2024 végéig
\end{enumerate}

\bigskip

%--- 1. feladat: árfolyam ---
\subsection{1. feladat -- Árfolyam grafikon}

\noindent Vizualizáljuk, hogyan alakult az Apple részvény záróárfolyama USD-ben
az elmúlt években. A \texttt{df} DataFrame \texttt{Close} oszlopa tartalmazza
a napi záróárakat.

\bigskip

\lbl{Feladat:}
Fogalmazz meg egy promptot a Gemini számára, amely megjeleníti a részvény árfolyamát!

\begin{prompt}[Javasolt prompt]
Te egy Python adatvizualizációs szakértő vagy.
### Feladat
Írj Python kódot, amely megjeleníti a df DataFrame Close oszlopát
idősoros vonaldiagramként.
### Formátum
- Legyen cím, tengelyfeliratok (x: Dátum, y: Árfolyam USD)
- figsize=(14, 5)
### Kontextus
df egy pandas DataFrame, dátum indexszel és Close oszloppal.
\end{prompt}

\noindent A kapott kódot másold be a notebook megfelelő cellájába és futtasd le.

\medskip\lbl{Módosítás -- jelöld meg a fontos időszakokat!}

\begin{prompt}[Módosító prompt]
Módosítsd az előző kódot úgy, hogy a következő időszakok
explicit látszódjanak a diagramon színezett háttérrel:
- COVID-összeomlás: 2020-02 -- 2020-06
- Ukrajna háború + kamatemelés: 2022-02 -- 2022-12
Minden időszakhoz legyen felirat is a diagramon.
\end{prompt}

\note{Próbáld ki más részvénnyel is! Cseréld le a \texttt{TICKER} értékét
(pl.\ \texttt{"MSFT"}, \texttt{"TSLA"}, \texttt{"NVDA"}) és futtasd újra
az adatletöltő cellát.}

\bigskip

%--- 2. feladat: napi hozam ---
\subsection{2. feladat -- Napi hozam}

\noindent Az árfolyam önmagában nem mutatja meg, mekkora volt a napi változás.
A napi hozam (daily return) ezt fejezi ki százalékban -- megmutatja,
mikor voltak a nagy esések és emelkedések.

\bigskip

\lbl{Feladat:}
Az előző prompt folytatásaként kérd meg a Geminit a napi hozam
kiszámítására és megjelenítésére!

\begin{prompt}[Javasolt prompt]
Az előző kód folytatásaként számítsd ki a napi százalékos hozamot
a Close oszlopból, és jelenítsd meg ugyanolyan stílusban.
Legyen cím és tengelyfeliratok.
\end{prompt}

\medskip\lbl{Módosítás -- statisztikák kiemelése!}

\begin{prompt}[Módosító prompt]
Módosítsd az előző kódot úgy, hogy a diagramon látszódjon:
- A legjobb és legrosszabb nap (pont + felirat a dátumal és értékkel)
- Egy vízszintes vonal a napi hozamok átlagánál
\end{prompt}

\bigskip

%--- 3. feladat: kumulált hozam ---
\subsection{3. feladat -- Kumulált hozam}

\noindent A kumulált hozam megmutatja: ha 2020 elején 1000 dollárt fektettél volna be,
ma mennyi lenne? Ezt a részvényre és az S\&P~500 indexre egyszerre vizualizáljuk --
így azonnal látszik, túlteljesítette-e a részvény a piacot.

\begin{itemize}[noitemsep, topsep=4pt]
  \item Ha a részvény vonala az index felett van $\rightarrow$ jobban teljesített a piacnál
  \item Ha alatta $\rightarrow$ egyszerűbb lett volna indexalapba fektetni
\end{itemize}

\bigskip

\lbl{Feladat:}
Fogalmazz meg egy promptot a Gemini számára mindkét eszköz kumulált hozamának
megjelenítéséhez!

\begin{prompt}[Javasolt prompt]
Az előző kód folytatásaként számítsd ki és jelenítsd meg a df és a df_index
kumulált hozamát egy diagramon, 1000 dolláros kezdőértékkel.
Legyen cím, tengelyfeliratok (y: Portfólió értéke USD) és jelmagyarázat.
\end{prompt}

\note{Ha valaki szeretné tudni a pontos végösszeget: kérd meg a Geminit, hogy
írja ki szövegesen is -- "1000 dollárból X lett a részvénnyel, Y az indexszel."}

\bigskip\noindent\rule{\linewidth}{0.4pt}\bigskip

%=============================================================
\section{2. RÉSZ -- FinBERT sentiment elemzés}
%=============================================================

\noindent\textit{Cél: Egy pénzügyi szövegekre finomhangolt nyelvi modell (FinBERT)
segítségével megvizsgáljuk, hogy az Apple-hez kapcsolódó hírek pozitív,
negatív vagy semleges hangulatot közvetítenek-e -- majd összevetjük az árfolyammal.}

\bigskip

\subsection{0. lépés -- Telepítés és modell betöltése}

\noindent Futtasd le sorban a notebook következő celláit:

\begin{enumerate}[noitemsep, topsep=4pt]
  \item \textbf{Telepítés / modell betöltés} -- a FinBERT modell letöltése és betöltése
  \item \textbf{Hírek betöltése} -- az \texttt{AAPL\_hirek.csv} fájl beolvasása
\end{enumerate}

\warn{A FinBERT modell (~400 MB) az első futtatáskor letöltodik a Hugging Face szervereiről.
Ez 1--2 percet vehet igénybe -- ez normális, csak egyszer kell letölteni!
A \texttt{AAPL\_hirek.csv} fájlt töltsd fel a Colab session-be a bal oldali fájlkezelőn keresztül.}

\bigskip

\subsection{1. feladat -- Sentiment elemzés futtatása}

\noindent A FinBERT minden hírhez három értéket ad: \textit{positive},
\textit{negative}, \textit{neutral} -- és mindegyikhez egy 0--1 közötti valószínűségi score-t.
A domináns sentiment az, amelyik a legnagyobb értékkel bír.

\bigskip

\lbl{Feladat:}
Fogalmazz meg egy promptot a Gemini számára, amely lefuttatja
a FinBERT elemzést az összes hírre!

\begin{prompt}[Javasolt prompt]
Te egy Python NLP szakértő vagy.
### Feladat
Írj kódot, amely a FinBERT modellel meghatározza minden hír
sentiment-jét és score-ját.
### Formátum
Az eredményt add hozzá a hirek DataFrame-hez:
- sentiment oszlop: "positive", "negative" vagy "neutral"
- score oszlop: a domináns sentiment valószínűsége (0-1)
### Kontextus
- hirek egy pandas DataFrame, date és headline oszlopokkal
- tokenizer és model már be van töltve (ProsusAI/finbert)
- A modell kimenetének sorrendje: positive, negative, neutral
\end{prompt}

\bigskip

\subsection{2. feladat -- Eredmények táblázatban}

\noindent Ha megvan a sentiment minden hírhez, jelenítsük meg
áttekinthető formában -- helyesen azonosítja-e a FinBERT
a pozitív és negatív híreket?

\bigskip

\lbl{Feladat:}

\begin{prompt}[Javasolt prompt]
Az előző kód folytatásaként jelenítsd meg a hirek DataFrame-et
táblázatban: dátum, hír címe, sentiment és score oszlopokkal.
A negatív sorokat emeld ki pirossal, a pozitívokat zölddel.
\end{prompt}

\tip{Ellenőrizd: a COVID-hírek (2020 február--március) negatívak,
a rekord bevételi bejelentések pozitívak? Ha igen, a modell jól működik.}

\bigskip

\subsection{3. feladat -- Top 5 legpozitívabb és legnegatívabb hír}

\noindent Melyik hír kapta a legmagasabb pozitív score-t?
Melyik a legnegatívabbat? Ez a két lista sokat elárul arról,
hogyan látja a modell az Apple híreit.

\bigskip

\lbl{Feladat:}

\begin{prompt}[Javasolt prompt]
Az előző kód folytatásaként jelenítsd meg a top 5 legpozitívabb
és top 5 legnegatívabb hírt két külön táblázatban, score szerint rendezve.
\end{prompt}

\bigskip

\subsection{4. feladat -- Sentiment megoszlás}

\noindent Összességében hogyan oszlanak meg a hírek?
Több pozitív vagy negatív hír jelent meg az Apple-ről 2020--2024 között?

\bigskip

\lbl{Feladat:}

\begin{prompt}[Javasolt prompt]
Az előző kód folytatásaként készíts kördiagramot
a sentiment kategóriák (positive/negative/neutral) megoszlásáról.
\end{prompt}

\bigskip

\subsection{5. feladat -- Hírek az árfolyamon}

\noindent Az utolsó vizualizáció: minden hír megjelenik az árfolyam
diagramon egy pontként -- zölddel a pozitív, pirossal a negatív híreket jelölve.
Vizuálisan látszik-e hogy a negatív hírek esések közelében vannak?

\warn{55 hírből statisztikai korrelációt számolni nem lenne megbízható.
Amit viszont látunk: a FinBERT helyesen azonosítja a negatív és pozitív híreket --
ez önmagában is értékes eredmény. Megbízható korreláció számításhoz
több száz hír kellene egyenletes időbeli elosztással.}

\bigskip

\lbl{Feladat:}

\begin{prompt}[Javasolt prompt]
Az előző kód folytatásaként jelenítsd meg az AAPL árfolyamát (df Close oszlop)
és rajzold rá a híreket pontként az árfolyam értékén:
zöld pont = positive, piros pont = negative, szürke pont = neutral.
Legyen jelmagyarázat és cím.
\end{prompt}

\note{Próbáld meg értelmezni együtt a diagramot: a 2020 februári
negatív hírek az árfolyamesés előtt vagy után jelennek meg?
Mit látunk 2023-ban a Berkshire-eladások idején?}

\bigskip\noindent\rule{\linewidth}{0.4pt}\bigskip

%=============================================================
\section{3. RÉSZ -- Kereskedési stratégia}
%=============================================================

\noindent\textit{Cél: Mozgóátlag alapú kereskedési stratégiát építünk,
backtesteljük a 2020--2024-es időszakra, és összehasonlítjuk
az egyszerű buy \& hold megközelítéssel.}

\bigskip

\subsection{0. lépés -- Technikai indikátorok}

\noindent Futtasd le a notebook következő celláját -- ez kiszámítja
az 50 és 200 napos mozgóátlagot és megjeleníti az árfolyamon.

\medskip
\noindent Az \textbf{50 napos mozgóátlag} a rövid távú trendet mutatja,
a \textbf{200 napos} a hosszú távú irányt. Ha a rövid átlag átlép
a hosszú fölé, az emelkedő trendet jelez -- és fordítva.

\bigskip

\subsection{1. feladat -- Kereskedési jelzések}

\noindent A két mozgóátlag keresztezési pontjai adják a kereskedési jelzéseket:

\begin{itemize}[noitemsep, topsep=4pt]
  \item \textbf{Golden cross} -- az MA50 átlép az MA200 fölé $\rightarrow$ vételi jel
  \item \textbf{Death cross} -- az MA50 átlép az MA200 alá $\rightarrow$ eladási jel
\end{itemize}

\bigskip

\lbl{Feladat:}
Fogalmazz meg egy promptot a Gemini számára, amely meghatározza
a keresztezési pontokat és megjeleníti őket az árfolyamon!

\begin{prompt}[Javasolt prompt]
Az előző kód folytatásaként azonosítsd a golden cross és death cross
pontokat a df DataFrame MA50 és MA200 oszlopai alapján.
Jelenítsd meg az árfolyamot a mozgóátlagokkal együtt,
és jelöld a golden cross pontokat zöld felfelé mutató háromszöggel,
a death cross pontokat piros lefelé mutató háromszöggel.
\end{prompt}

\bigskip

\subsection{2. feladat -- Backtesting}

\noindent Ha 2020 elején 1000 dollárral indultunk és minden jelzésnél
vettünk/eladtunk, mennyit ér most a portfóliónk?

\bigskip

\lbl{Feladat:}

\begin{prompt}[Javasolt prompt]
Az előző kód folytatásaként szimulálj egy egyszerű backtestet
1000 dolláros kezdőtőkével:
- Golden cross-nál: teljes tőkével vétel
- Death cross-nál: teljes pozíció eladása
Számítsd ki és írd ki a végső portfólió értékét.
\end{prompt}

\tip{Ha a stratégia rosszabbul teljesített mint a buy \& hold,
ez is értékes eredmény -- megmutatja hogy a mozgóátlag stratégia
nem mindig veri meg a piacot. Ez a valóságban is így van!}

\bigskip

\subsection{3. feladat -- Stratégia vs. Buy \& Hold}

\noindent Az utolsó vizualizáció: hogyan alakult a portfólió értéke időben --
a mozgóátlag stratégiával és az egyszerű buy \& hold megközelítéssel.

\bigskip

\lbl{Feladat:}

\begin{prompt}[Javasolt prompt]
Az előző kód folytatásaként jelenítsd meg egy diagramon
a mozgóátlag stratégia és a buy \& hold portfólió értékét időben,
1000 dolláros kezdőtőkéről indulva.
Legyen cím, tengelyfeliratok (y: Portfólió értéke USD) és jelmagyarázat.
\end{prompt}

\note{Gondolkozzatok együtt: mikor teljesít jól a mozgóátlag stratégia?
Erősen trendelő piacokon igen -- oldalazó vagy hirtelen forduló piacokon kevésbé.
Az Apple 2020--2024 között erősen emelkedő trend volt -- ez kedvez a buy \& hold-nak.}

\bigskip

\subsection{4. feladat -- Módosítás: más paraméterek}

\noindent A stratégia eredménye erősen függ a választott paraméterektől.
Mi történik ha rövidebb mozgóátlagokat használunk?

\bigskip

\lbl{Módosító prompt:}

\begin{prompt}[Módosító prompt]
Módosítsd az előző kódot úgy, hogy MA20 és MA50 paramétereket
használjon MA50 és MA200 helyett.
Futtasd újra a backtestet és írd ki az új végső portfólió értékét.
\end{prompt}

\note{Rövidebb paraméterek = több jelzés, gyorsabb reakció, de több hamis jelzés is.
Hosszabb paraméterek = kevesebb jelzés, lassabb reakció, de megbízhatóbb trend.}

\bigskip\noindent\rule{\linewidth}{0.4pt}\bigskip

%=============================================================
\newpage
\section*{Cheat Sheet -- Vibe coding alapok}
%=============================================================

\begin{tcolorbox}[cheatstyle, title={A vibe coding lépései}]
\begin{tabularx}{\linewidth}{@{}lX@{}}
\textbf{1. Értsd meg a feladatot} & Mi az input? Mi az elvárt output? \\[3pt]
\textbf{2. Írd meg a promptot} & RTFC módszertan -- szerep, feladat, formátum, kontextus \\[3pt]
\textbf{3. Futtasd a kódot} & Másold be a Colab cellába, futtasd le \\[3pt]
\textbf{4. Értékeld az eredményt} & Helyes? Érthető? Szép? \\[3pt]
\textbf{5. Iterálj} & Ha nem jó, pontosítsd a promptot -- ne a kódot \\
\end{tabularx}
\end{tcolorbox}

\smallskip

\begin{tcolorbox}[cheatstyle, title={Hasznos prompt kiegészítések kódgeneráláshoz}]
\begin{tabularx}{\linewidth}{@{}lX@{}}
\textit{"Magyarázd el a kódot röviden."} & Értsd meg amit futtatsz \\[3pt]
\textit{"Adj hozzá kommenteket."} & Olvashatóbb kód \\[3pt]
\textit{"A df változó neve ne változzon."} & Konzisztencia a notebookban \\[3pt]
\textit{"Ne importálj új csomagokat."} & Ha már telepítve van minden \\[3pt]
\textit{"Hibaüzenetet kaptam: [ide másold be]. Javítsd ki."} & Hibakeresés \\
\end{tabularx}
\end{tcolorbox}

\smallskip

\begin{tcolorbox}[cheatstyle, title={Notebook paraméterek -- gyors változtatás}]
\begin{tabular}{@{}ll@{}}
\texttt{TICKER = "MSFT"}        & Microsoft \\[2pt]
\texttt{TICKER = "TSLA"}        & Tesla \\[2pt]
\texttt{TICKER = "NVDA"}        & NVIDIA \\[2pt]
\texttt{TICKER = "\^{}GSPC"}    & S\&P 500 index \\[2pt]
\texttt{KEZDO\_DATUM = "2022-01-01"} & Csak a legutóbbi 3 év \\
\end{tabular}
\end{tcolorbox}

\smallskip

\begin{tcolorbox}[warnstyle]
\textbf{\textcolor{warnborder}{Fontos:}}
Ez az elemzés kizárólag oktatási célokat szolgál.
A modellek által generált előrejelzések és kereskedési jelzések
nem minősülnek befektetési tanácsnak.
A valós befektetési döntésekhez mindig konzultálj szakértővel.
\end{tcolorbox}

\vfill
\begin{center}
\small\textcolor{medgray}{Sapientia EMTE \textbullet{} Generatív AI és üzleti alkalmazások \textbullet{} 2026}
\end{center}

\end{document} 